% !TeX program = uplatex
% !TeX encoding = UTF-8

\documentclass[uplatex,dvipdfmx,11pt]{bxjsarticle}

% ====== ページ設定（bxjsarticle既定。警告回避のため余白は必要時のみ設定） ======
% 余白を変更したい場合は次の行を有効化してください。
% \setpagelayout{margin=25mm}

% ====== 数式・記号 ======
\usepackage{amsmath,amssymb,mathtools}
\usepackage{bm}

% ====== 図表 ======
\usepackage{graphicx}
\usepackage{subcaption}
\usepackage{booktabs}
\usepackage{siunitx}

% ====== 参照・リンク ======
\usepackage[dvipdfmx]{hyperref}
\usepackage{pxjahyper}
\usepackage[nameinlink,noabbrev]{cleveref}

% ====== 箇条書き等 ======
\usepackage{enumitem}

% ====== 文献（biblatex + biber） ======
\usepackage[backend=biber,style=numeric,sorting=none]{biblatex}
% 本テンプレートは tex/ 配下にあるため、tex/ 直下の references.bib を参照
\addbibresource{../references.bib}

% ====== 自作コマンド（必要に応じて拡張） ======
\newcommand{\SF}{\textsf{SF}}
\newcommand{\ER}{\textsf{ER}}
\newcommand{\RR}{\textsf{RR}}

\title{（タイトル）スケールフリー上の次数依存SAR解析テンプレート}
\author{（氏名・所属）}
\date{\today}

\begin{document}
\maketitle

\begin{abstract}
概要を1段落で。目的・方法・主な所見を簡潔に記述します。
\end{abstract}

\section{背景と目的}
関連研究や位置付けを簡潔に。例えば次数依存パーコレーションの条件\cite{BaxterTimar2021}や、
SARにおける不連続転移の報告\cite{WangEtAl2015} などを示します。

\section{モデルと実装}
\subsection{ネットワーク}
\begin{itemize}[leftmargin=2em]
  \item 構成モデル（\texttt{Config}）：$P(k)\propto k^{-\gamma}$，$k_{\min}$，$\langle k\rangle$ など
  \item ER/BA/RR も利用可（必要に応じて切替）
\end{itemize}

\subsection{SARダイナミクス}
非冗長情報メモリ・閾値$T$を採用。伝播率は
\begin{equation}
\lambda\,k_u^{\alpha}k_v^{\beta}\ (\beta=0~\text{例})
\end{equation}
の形で次数依存性を与えます。

\section{設定と実行}
\subsection{代表設定（例）}
\begin{table}[htbp]
  \centering
  \caption{代表的な実験設定（編集して使用）}
  \begin{tabular}{@{}ll@{}}\toprule
    項目 & 値（例） \\ \midrule
    ネットワーク & Config：$\gamma=2.3$, $k_{\min}=5$, $\langle k\rangle=10$ \\
    ノード数 & $50{,}000$ \\
    閾値 $T$ & $1$，活動家割合 $p=0.0$ \\
    初期感染数 & $k_0=1$ \\
    回復率 & $\gamma=1.0$，観測時間 $t_{\max}=200$ \\
    掃引 & $\lambda\in[0,3],\ \Delta\lambda=0.01$，$\alpha\in\{-2.5,-2.0,-1.0,-0.5,0.0\}$ \\
    出力 & 終値（\texttt{isFinal=true}） \\
  \bottomrule\end{tabular}
\end{table}

\subsection{コマンド例}
\begin{verbatim}
./run-fastsar.sh
# もしくは
./gradlew :app:runFastSAR
\end{verbatim}
出力は \verb|out/fastsar/.../results_{00..}.csv| に追記されます。

\section{結果}
\subsection{\texorpdfstring{$R(\infty)$}{R(infty)} vs \texorpdfstring{$\lambda$}{lambda}}
\begin{figure}[htbp]
  \centering
  % \includegraphics[width=0.85\linewidth]{fig/Rinf_vs_lambda_example.pdf}
  \caption{最終規模と伝播率の関係（例）。}
\end{figure}

\subsection{表: 転移点推定とジャンプ量}
\begin{table}[htbp]
  \centering
  \begin{tabular}{@{}lllll@{}}\toprule
    ネットワーク & $\gamma$ & $\alpha$ & 推定 $\lambda_c$ & $\Delta R$ \\ \midrule
    SF & 2.5 & -0.5 & 0.32 & 0.41 \\
    ER & --  & 0.0  & 0.58 & 0.35 \\
  \bottomrule\end{tabular}
  \caption{転移点推定の例（数値はダミー）。}
\end{table}

\section{考察と今後の課題}
所見の要約，感度分析の計画，理論的検討の方向性など。

\printbibliography

\end{document}
